\documentclass[11pt]{amsart}
\usepackage{amsmath,amssymb,amsthm,mathtools}
\usepackage[margin=1in]{geometry}
\usepackage{hyperref}
\usepackage{array}
\usepackage{booktabs}
\usepackage{longtable}
\usepackage{xcolor}
\usepackage{enumitem}
\usepackage{tcolorbox}
\tcbuselibrary{breakable,skins}

%% Theorem environments
\newtheorem{theorem}{Theorem}[section]
\newtheorem{proposition}[theorem]{Proposition}
\newtheorem{lemma}[theorem]{Lemma}
\newtheorem{corollary}[theorem]{Corollary}
\newtheorem{conjecture}[theorem]{Conjecture}
\theoremstyle{definition}
\newtheorem{definition}[theorem]{Definition}
\newtheorem{example}[theorem]{Example}
\theoremstyle{remark}
\newtheorem{remark}[theorem]{Remark}
\newtheorem{notation}[theorem]{Notation}

%% Colored boxes for walls and no-gos
\newtcolorbox{wallbox}[1]{
  colback=red!5!white, colframe=red!60!black,
  title={\textbf{Wall #1}},
  breakable
}
\newtcolorbox{nogobox}[1]{
  colback=blue!5!white, colframe=blue!60!black,
  title={\textbf{No-Go #1}},
  breakable
}
\newtcolorbox{resultbox}[1]{
  colback=green!5!white, colframe=green!60!black,
  title={\textbf{Result: #1}},
  breakable
}

%% Math macros
\newcommand{\RH}{\mathrm{RH}}
\renewcommand{\Re}{\operatorname{Re}}
\newcommand{\half}{\tfrac{1}{2}}
\newcommand{\om}{\omega}
\newcommand{\kap}{\kappa}
\newcommand{\Fq}{F_q}
\newcommand{\Gq}{G_q}
\newcommand{\sig}{\sigma}
\newcommand{\gam}{\gamma}
\newcommand{\Lam}{\Lambda}
\newcommand{\eps}{\varepsilon}
\newcommand{\Mk}{M_k}
\newcommand{\Ck}{C_k}
\newcommand{\supp}{\operatorname{supp}}
\newcommand{\Spec}{\operatorname{Spec}}

\title[The Architecture of Obstruction]{The Architecture of Obstruction:\\ A Complete Mathematical Map of the Riemann Hypothesis Program from $\omega$-Class Statistics to the Symmetry/Positivity Wall}

\author{David Alejandro Trejo Pizzo}
\thanks{Independent Researcher. \texttt{dtrejopizzo@gmail.com}}


\begin{document}
\maketitle

\begin{abstract}
This paper is a comprehensive mathematical survey of the analytic, functional-analytic,
and geometric approaches to the Riemann Hypothesis developed by the author.
Unlike a standard survey, its primary aim is not to catalogue what is known, but to
\emph{map every identified obstruction with full mathematical precision}: to state exactly
which theorem, computation, or structural argument terminates each route, and why.

The program explored six distinct families of approaches:
(I) the $\om$-class program (statistics of $\om(n)$, the number of distinct prime factors);
(II) the localized Weil quadratic form (explicit-formula-based detector, Kreĭn index);
(III) the capstone paradigm (wrong-sign obstruction and its eight manifestations);
(IV) the geometric programs (Spec~$\mathbb{Z}$ cohomology, Lefschetz dichotomy, SURF);
(V) the functional-analytic arc (Kreĭn/Pontryagin spaces, de Branges axioms);
(VI) the information barrier (exact classification of which $\om$-class objects carry
zero information vs.\ RH-equivalent information).

We identify three surviving walls: \textbf{W1} (Weil positivity/PNT, the classical wall
in new dress), \textbf{W2} (de Branges canonical systems, requiring a global
Spec~$\mathbb{Z}$ geometry), and \textbf{W3} (Montgomery pair-correlation / GUE gap
universality).  Eight no-go theorems (N1--N8) are stated in full.  Twenty-seven
genuinely new results are listed that are \emph{independent of RH}.

The paper ends with a synthesis theorem: \emph{the reason RH resists all known
approaches is that the zero-location constraint $\kap=0$ is (i) global, (ii) unilateral
(an upper bound), and (iii) orthogonal to all available symmetry data.  Every programme
that terminates does so because it supplies only items (i-iii) partially.}  Nothing here
proves RH.

\medskip
\noindent\textbf{Key words and phrases.}
Riemann Hypothesis, explicit formula, Weil positivity, Kreĭn index, Pontryagin spaces,
de Branges canonical systems, $\om$-class, multiplicative chaos, information barrier,
Selberg--Delange, obstructions, no-go theorems, capstone, wrong-sign obstruction.
\end{abstract}

\tableofcontents
\newpage

\newpage

%% =====================================================================
\section*{Preface: on the use of this document}
%% =====================================================================

This is not a proof of the Riemann Hypothesis.  It is not a record of progress toward a
proof.  It is a precise, mathematically honest account of what a systematic multi-year
programme found: which directions work, which fail, and \emph{exactly why} each one
fails.

The programme operates under one non-negotiable principle: a false positive is worse
than a failure.  Every claim is either proved or explicitly flagged as conjecture.
Every refutation of a previous claim within the programme is documented.  The reader
will find several places where earlier results were corrected or retracted; these are
not embarrassments but structural data.

The paper is organized as a guided tour of the obstruction landscape, not a
chronological account.  The guiding question is always the same: \emph{at what point
does this route encounter an obstacle it cannot pass, and is that obstacle fundamental
or contingent?}

\bigskip
\noindent\textbf{How to read this paper.}
Each major section describes one family of approaches.  Within each section, the
mathematics is developed from scratch (basic definitions, key theorems, proofs or
proof sketches) before the wall is identified.  Readers familiar with the area can
skip to the ``Wall'' box at the end of each subsection.  The synthesis is in
\S\ref{sec:synthesis}.

%% =====================================================================
\section{Foundations and Framework}
\label{sec:foundations}
%% =====================================================================

\subsection{The Riemann Hypothesis and its equivalent forms}
\label{ss:RH}

Let $\zeta(s) = \sum_{n=1}^\infty n^{-s}$ for $\Re(s)>1$, extended meromorphically to
$\mathbb{C}$.  Its non-trivial zeros $\rho = \sig_\rho + i\gam_\rho$ satisfy
$0 < \sig_\rho < 1$ and come in pairs $\rho, 1-\bar\rho$.

\begin{definition}[The Riemann Hypothesis]
$\RH$: every non-trivial zero $\rho$ of $\zeta(s)$ satisfies $\sig_\rho = \half$.
\end{definition}

The programme uses four equivalent forms of RH.  We state each precisely.

\begin{proposition}[Equivalent forms of RH used in this programme]
\label{prop:RH-equiv}
The following are equivalent:
\begin{enumerate}[label=\upshape(\roman*)]
  \item \textup{(RH)} Every non-trivial zero of $\zeta$ has real part $\half$.
  \item \textup{(PNT-error)} For all $\eps>0$,
    $\psi(x) = x + O(x^{1/2+\eps})$, where $\psi(x)=\sum_{n\le x}\Lam(n)$.
  \item \textup{(Weil positivity)} The Weil explicit-formula form $W_f \ge 0$ for
    all admissible test functions $f$.
  \item \textup{(Kreĭn index)} The Kreĭn index $\kap(Q)$ of the localized Weil form
    $Q$ equals zero.
\end{enumerate}
\end{proposition}

\begin{proof}
(i)$\Leftrightarrow$(ii): classical Ingham/Landau.
(i)$\Leftrightarrow$(iii): Weil~\cite{Weil52}; see \S\ref{ss:Weil}.
(iii)$\Leftrightarrow$(iv): \cite{RP8}; see \S\ref{ss:Krein}.
\end{proof}

The localized Weil form $Q$ and its Kreĭn index $\kap$ are the central objects of the
programme; we describe them fully in \S\ref{sec:Weil}.

\subsection{The explicit formula}
\label{ss:explicit}

The most-used tool throughout the programme is the explicit formula.  Let
$\hat f(s) = \int_0^\infty f(x)x^{s-1}dx$ be the Mellin transform.

\begin{theorem}[Explicit formula]
\label{thm:explicit}
For an admissible test function $f$ (smooth, compactly supported away from 0 and
prime powers),
\begin{equation}
  \label{eq:explicit}
  \sum_{n=1}^\infty \Lam(n)f(n)
  = \hat f(1) - \sum_\rho \hat f(\rho) - \frac{\hat f(0)+\hat f(1)}{2}
  - \sum_{k=1}^\infty \hat f(-2k),
\end{equation}
where the sum $\sum_\rho$ is over non-trivial zeros, ordered by $|\Im(\rho)|$.
\end{theorem}

The term $\sum_\rho \hat f(\rho)$ is the zero contribution.  The programme repeatedly
encounters objects whose ``zero sensitivity'' is determined by whether this term is
present and non-trivially computable.

\subsection{The Selberg--Delange method}
\label{ss:SD}

The $\om$-class programme relies heavily on the following.

\begin{theorem}[Selberg--Delange, cf.\ Tenenbaum \cite{Ten15}~\S II.5]
\label{thm:SD}
Let $F(s) = \sum_{n=1}^\infty a(n)n^{-s} = \zeta(s)^z G(s)$ where $z\in\mathbb{C}$
and $G$ is holomorphic and non-vanishing in $\Re(s)\ge\half$.  Then as $x\to\infty$,
\[
  \sum_{n\le x} a(n) \;=\;
  \frac{G(1)}{\Gamma(z)} x (\log x)^{z-1}
  \left[ 1 + \sum_{k=1}^K \frac{c_k}{\log^k x} + O\!\left(\frac{1}{\log^{K+1} x}\right)\right],
\]
where the constants $c_k$ are determined by the Laurent expansion of $G$ and $\zeta$
at $s=1$.  No non-trivial zero of $\zeta$ appears in this expansion at any finite order.
\end{theorem}

\begin{remark}
The crucial point for the programme: zeros of $\zeta$ do not appear in the
Selberg--Delange expansion because $\zeta^z$ has a \emph{zero} (of order $z$) at each
$\rho$, never a pole.  This ``zeros-are-zeros'' observation is the foundation of the
information barrier; see \S\ref{sec:omega-full}.
\end{remark}

\subsection{Notation summary}
\label{ss:notation}

Throughout the paper:
\begin{itemize}
\item $\om(n)$ = number of \emph{distinct} prime factors of $n$;
  $\Omega(n)$ = total count with multiplicity.
\item $q = k^2$ for a parameter $k > 0$.
\item $F_q(s) = \sum_{n\ge 1} q^{\om(n)} n^{-s} = \zeta(s)^q G_q(s)$, where
  $G_q$ is described in \S\ref{sec:omega-full}.
\item $M_k(N) = \sum_{n\le N} q^{\om(n)}/n$, $C_k(h,N) = \sum_{n\le N} q^{\om(n)}q^{\om(n+h)}/n$.
\item $\kap = \kap(Q)$ = Kreĭn index of the localized Weil form; equals the number
  of off-line zeros (under the programme's working hypothesis).
\item $\sig_0 = \sup\{\Re\rho : \zeta(\rho)=0,\; \rho \text{ non-trivial}\}$.
  RH says $\sig_0 = \half$.
\end{itemize}

%% =====================================================================
\section{The $\omega$-Class Programme: from Statistics to Information Barrier}
\label{sec:omega-full}
%% =====================================================================

\subsection{Motivation: moment statistics of $\zeta$ on the critical line}
\label{ss:omega-motivation}

The Keating--Snaith conjecture~\cite{KS00} predicts that the moments of $|\zeta|$ on
the critical line satisfy
\[
  \frac{1}{T}\int_0^T |\zeta(\half + it)|^{2k} dt \sim C_k (\log T)^{k^2}.
\]
The exponent $k^2$ is remarkable: it is the second cumulant of a log-correlated field.
The programme's starting point: can the exponent $k^2$ be derived from the statistics
of $\om(n)$, \emph{without} a priori knowledge of the zeros?

\subsection{The Dirichlet series of $q^{\omega(n)}$}
\label{ss:dirichlet}

\begin{proposition}[Factorization]
\label{prop:factorization}
For $q > 0$ and $\Re(s) > 1$:
\[
  F_q(s) = \sum_{n=1}^\infty \frac{q^{\om(n)}}{n^s}
  = \prod_p \frac{1+(q-1)p^{-s}}{1-p^{-s}}
  = \zeta(s)^q \cdot G_q(s),
\]
where
\[
  G_q(s) = \prod_p \bigl[(1+(q-1)p^{-s})(1-p^{-s})^q\bigr].
\]
\end{proposition}

\begin{proof}
At each prime $p$, $\om(p^k) = 1$ for all $k\ge 1$, so the local Euler factor is
\[
  \sum_{k=0}^\infty \frac{q^{\om(p^k)}}{p^{ks}}
  = 1 + \sum_{k=1}^\infty \frac{q}{p^{ks}}
  = 1 + \frac{q}{p^s-1}
  = \frac{1+(q-1)p^{-s}}{1-p^{-s}}.
\]
The factorization $= \zeta(s)^q G_q(s)$ follows from
$\zeta(s)^q = \prod_p (1-p^{-s})^{-q}$, giving
$G_q(s) = \prod_p [(1+(q-1)p^{-s})(1-p^{-s})^q]$.
\end{proof}

\begin{lemma}[L1: holomorphy of $G_q$]
\label{lem:Gq}
For every $q > 0$, $G_q(s)$ is holomorphic and non-vanishing in $\Re(s) > \half$.
\end{lemma}

\begin{proof}
Write $\log G_q(s) = \sum_p \log[(1+(q-1)p^{-s})(1-p^{-s})^q]$.
For $\sig = \Re(s) > \half$:
\begin{align*}
\log[(1+(q-1)p^{-s})(1-p^{-s})^q]
&= (q-1)p^{-s} - qp^{-s} + O(p^{-2\sig}) \\
&= -p^{-s} + O(p^{-2\sig}).
\end{align*}
The sum $\sum_p p^{-s}$ converges absolutely for $\sig > 1$, and for $\sig > \half$
the tail from the $O(p^{-2\sig})$ terms converges.  The difference $\sum_p p^{-s} -
\log\zeta(s)$ converges absolutely for $\sig > \half$ (Mertens estimate), giving
holomorphy of $\log G_q$ in this half-plane, hence holomorphy and non-vanishing of $G_q$.
\end{proof}

\subsection{The key structural consequence: zeros of $\zeta$ are zeros of $F_q$}
\label{ss:zerosofzeros}

\begin{theorem}[Structural theorem on $M_k(N)$]
\label{thm:Mk-monotone}
For every $q > 0$:
\begin{enumerate}[label=\upshape(\roman*)]
  \item $F_q(\rho) = 0$ for every non-trivial zero $\rho$ of $\zeta$.
  \item The Perron integral for $M_k(N) = \sum_{n\le N} q^{\om(n)}/n$ collects no
    residue at any non-trivial zero of $\zeta$.
  \item $M_k(N)$ is strictly increasing in $N$ and has the asymptotic expansion
    \[
      M_k(N) = \frac{G_q(1)}{\Gamma(q)}(\log N)^q
      + \sum_{j=1}^{J} c_j (\log N)^{q-j} + O((\log N)^{q-J-1})
    \]
    with computable arithmetic coefficients $c_j$ and no oscillatory terms $N^{\rho-1}$.
\end{enumerate}
\end{theorem}

\begin{proof}
(i): $F_q(\rho) = \zeta(\rho)^q G_q(\rho)$.  Since $\zeta(\rho)=0$ and $G_q(\rho)\ne 0$
(Lemma~\ref{lem:Gq}), we have $F_q(\rho)=0$.

(ii): By Perron's formula,
$M_k(N) = (2\pi i)^{-1}\int_{c-i\infty}^{c+i\infty} F_q(s)\frac{N^s}{s}\,ds$
for $c>1$.  Moving the contour left, residues arise at poles of $F_q(s)/s$.  Since
$F_q(\rho)=0$, the integrand has no pole at $\rho$; the only pole is at $s=1$
(of order $q$) and at $s=0$.

(iii): The Laurent expansion at $s=1$ gives (ii) of the Selberg--Delange theorem.
Each term $q^{\om(n)}/n > 0$ ensures strict monotonicity.
\end{proof}

\begin{remark}[The monotonicity barrier]
Theorem~\ref{thm:Mk-monotone}(iii) destroys any hope of detecting RH by measuring
``oscillations of $M_k(N)$'': there are none.  The function $M_k$ is completely
determined by the singularity of $F_q$ at $s=1$, which is arithmetic (depends only
on $G_q(1)$ and the Laurent coefficients), not analytic (no zero contributions).

This was a genuine surprise in the programme.  Phase~20 of the investigation began
with the hypothesis that $M_k(N)/(\log N)^{k^2}$ would oscillate with amplitude
$N^{-1/2}$ if RH holds and $N^{\sig_0-1}$ if a zero at $\sig_0>\half$ exists.
Theorem~\ref{thm:Mk-monotone} refutes this completely.
\end{remark}

\subsection{The k² exponent: derivation from Erd\H{o}s--Kac}
\label{ss:k2}

A central result of Phase~19~\cite{RP9} is that the Keating--Snaith exponent $k^2$
is a \emph{Poisson invariant} of the distribution of $\om(n)$, with no reference to
zeros of $\zeta$.

\begin{theorem}[T2: $k^2$ from Erd\H{o}s--Kac]
\label{thm:k2}
Let $\Mk(N) = \sum_{n\le N} q^{\om(n)}/n$.  The constant $q = k^2$ in the asymptotic
$\Mk(N) \sim C_q (\log N)^q$ is determined purely by the Poisson distribution of
$\om(n)$:
\[
  k^2 = \lim_{N\to\infty} \frac{\log \Mk(N)}{\log\log N}
  = \left.\frac{d^2}{dt^2}\right|_{t=0} \Lambda_\om(t),
\]
where $\Lambda_\om(t) = \log\mathbb{E}[e^{t\om(n)}] = (e^t-1)\log\log N + O(1)$
is the log-moment-generating function.  Explicitly, $k^2 = (e^t-1)\big|_{t=\log k}$
recovers $k^2$ from the inverse Legendre transform.

The large-deviation rate function is
\[
  I(\alpha) = \alpha\log\alpha - \alpha + 1 \quad(\text{Poisson with mean } 1),
\]
and corrections to the leading Poisson asymptotics are $O(1/\log\log N)$ with
Mertens-type arithmetic coefficients.  No non-trivial zero of $\zeta$ appears.
\end{theorem}

\begin{proof}[Proof sketch]
The Erd\H{o}s--Kac theorem gives $(\om(n)-\log\log n)/\sqrt{\log\log n}\to\mathcal{N}(0,1)$
in distribution.  In the large-deviation regime, $\om(n)\approx\alpha\log\log N$, the
log-MGF is $\Lambda_\om(t) = \sum_p p^{-1}(e^t-1)/(1+O(p^{-1}))\approx(e^t-1)\log\log N$.
The G\"artner--Ellis theorem gives $I(\alpha) = \sup_t\{\alpha t - (e^t-1)\} =
\alpha\log\alpha - \alpha + 1$.  The exponent $k^2 = z$ in $F_q = \zeta^z G_z$
satisfies $z = e^{\log k^2/\log\log N \cdot \log\log N}\approx k^2$ in the relevant
scaling; see~\cite{RP9} for the precise derivation.
\end{proof}

\subsection{The shifted correlation and the information barrier}
\label{ss:correlation}

The multiplicative function $q^{\om(n)}$ is insensitive to the zeros of $\zeta$.
A natural question: does \emph{any} statistic of the $\om$-class see the zeros?

\begin{definition}[Shifted correlation]
For $h \ge 1$ fixed and $q = k^2 > 0$:
\begin{align}
  C_k(h,N) &= \sum_{n\le N}\frac{q^{\om(n)}\,q^{\om(n+h)}}{n}, \\
  r_q(h) &= \lim_{N\to\infty}\frac{C_k(h,N)}{M_k(N)^2}
           = \prod_{p\mid h} f_q(p),
\end{align}
where $f_q(p) = p(1+(q-1)/p)^2/(p-1+q)$.
\end{definition}

The limit $r_q(h)$ is an arithmetic constant (finite Euler product over $p\mid h$),
independent of the zeros of $\zeta$.  The \emph{rate of convergence} to this limit,
however, encodes RH.

\begin{theorem}[T4: explicit error formula]
\label{thm:T4}
Fix $h\ge 1$ and $q > 0$.  Using the inclusion-exclusion identity
$q^{\om(n)} = \sum_{d\mid n,\,d\,\mathrm{sqfree}}(q-1)^{\om(d)}$ twice:
\[
  C_k(h,N) = S_q(h)\log N + K_q(h)
  - \frac{A_q(h)}{C_q^2}\sum_\rho \frac{N^{\rho-1}}{\rho}(\log N)^{2q-1}
  + O\!\bigl((\log N)^{2q-2}\bigr),
\]
where $S_q(h)$, $K_q(h)$ are arithmetic constants, $C_q = G_q(1)/\Gamma(q)$,
and the arithmetic coefficient of the zero sum is
\begin{equation}
  \label{eq:Aq}
  A_q(h) = \prod_{p\nmid h}\!\!\Bigl(1+\frac{2(q-1)}{p}\Bigr)
            \cdot\prod_{p\mid h}\!\!\Bigl(1+\frac{2(q-1)}{p}+\frac{(q-1)^2}{p}\Bigr)
  \;\; > \;\; 0.
\end{equation}
Consequently:
\[
  \frac{C_k(h,N)}{M_k(N)^2} - r_q(h)
  \;=\; -\frac{A_q(h)}{C_q^2}\sum_\rho \frac{N^{\rho-1}}{\rho} + O\!\left(\frac{1}{\log N}\right).
\]
\end{theorem}

\begin{proof}
Each inner sum $\sum_{n\le N:\,d\mid n,\,e\mid(n+h)}1/n$ (over an arithmetic progression
mod $\operatorname{lcm}(d,e)$) has error $-\phi(m)^{-1}\sum_\chi\bar\chi(r)\sum_{\rho_\chi}
N^{\rho_\chi-1}/\rho_\chi$ by the Dirichlet character/Perron decomposition.  For the
principal character $\chi_0$, this contributes zeros of $\zeta$.  Summing over all
squarefree pairs $(d,e)$ with $\gcd(d,e)\mid h$ and computing the Euler product of
local factors gives~\eqref{eq:Aq}.  Each local factor is a positive real number
(both numerator factors are positive for $q>0$, $p\ge 2$), so $A_q(h)>0$.
\end{proof}

\begin{theorem}[T4.1: RH-equivalence]
\label{thm:T4.1}
$\RH \Longleftrightarrow$ for every fixed $h\ge 1$ and $q > 0$:
\[
  \Bigl|\frac{C_k(h,N)}{M_k(N)^2} - r_q(h)\Bigr| = O(N^{-1/2+\eps}).
\]
\end{theorem}

\begin{proof}
Under RH all $\rho$ satisfy $|N^{\rho-1}| = N^{-1/2}$, and the zero-counting bound
gives the $O(N^{-1/2+\eps})$ estimate.  Conversely, an off-line zero $\rho_0 =
\sig_0+i\gam_0$ with $\sig_0>\half$ contributes $|A_q(h)/C_q^2| \cdot N^{\sig_0-1} /
|\rho_0| \gg N^{\sig_0-1} \to \infty$ faster than $N^{-1/2}$, contradicting
$O(N^{-1/2+\eps})$.
\end{proof}

\begin{theorem}[T5: impossibility of algebraic cancellation]
\label{thm:T5}
For any positive weights $q_1,\ldots,q_r>0$ and shifts $h_1,\ldots,h_r\ge 0$, the
arithmetic coefficient $A_{q_1,\ldots,q_r}(h_1,\ldots,h_r)$ multiplying
$\sum_\rho N^{\rho-1}/\rho$ in the error of the multi-point correlation
$D_r(N) = \sum_{n\le N}\prod_{j=1}^r q_j^{\om(n+h_j)}/n$ satisfies
$A_{q_1,\ldots,q_r}(h_1,\ldots,h_r) > 0$ strictly.
\end{theorem}

\begin{proof}
$A = \prod_p[\text{local factor at }p]$.  Each local factor is a weighted sum of
products $\prod_j q_j^{[\text{event}]}$ over residue classes, with all $q_j>0$.
Hence all summands are positive, all local factors are positive, and $A>0$.
\end{proof}

\begin{remark}[Duality of positivity]
Theorems~\ref{thm:Mk-monotone} and~\ref{thm:T5} are two faces of the same coin.
The positivity $q^{\om(n)}\ge 1$ guarantees (a) monotonicity of $M_k$ (no zero
oscillations; this is helpful), and (b) positivity of all Euler-product coefficients
(no algebraic cancellation; this is an obstruction).  No combination of positive-weight
$\om$-class correlations can eliminate the zero contribution.
\end{remark}

\subsection{Complete classification table}

Table~\ref{tab:omega-class} gives the complete classification of natural $\om$-class objects.

\begin{table}[h]
\centering
\renewcommand{\arraystretch}{1.4}
\begin{tabular}{lllll}
\toprule
Object & Asymptotics & Rate & Zeros & Class \\
\midrule
$M_k(N)/(\log N)^q \to C_q$ & $C_q$ & $O(1/\log N)$ & No (Laurent) & B \\
$M_k(N)/M_1(N)^q \to C_q$ & $C_q$ & $O(1/\log N)$ & No (cancel.) & B \\
$r_q(h)=\lim C_k/M_k^2$ & Euler prod. & exists & No (limit) & B \\
$I(\alpha)=\alpha\log\alpha-\alpha+1$ & Poisson & --- & No & B \\
$k^2 = $ exponent of $M_k$ & arithmetic & --- & No & B \\
$C_k(h,N)/M_k^2 - r_q(h)$ & 0 & $O(N^{-1/2+\eps})|$RH & Yes ($\zeta$) & A \\
$\Psi_q(x) = \sum q^{\om(n)}\Lam(n)$ & $q\cdot\psi(x)$ & $O(x^{-1/2+\eps})|$RH & Yes & C \\
\bottomrule
\end{tabular}
\caption{Information-barrier classification. Class A = RH-equivalent; B = new
math, RH-independent; C = reduces trivially to known objects.}
\label{tab:omega-class}
\end{table}

\begin{wallbox}{W1 (from the $\omega$-class direction)}
\textbf{Theorem T4.1} shows that the only RH-sensitive window in the $\om$-class is
the convergence rate of $C_k(h,N)/M_k(N)^2$.  Proving that rate is $O(N^{-1/2+\eps})$
is \emph{equivalent} to RH, via the PNT in arithmetic progressions.  This is the
classical wall W1 in new dress.  \textbf{Theorem T5} shows that no algebraic
cancellation within the $\om$-class can bypass this wall.
\end{wallbox}

%% =====================================================================
\section{The Localized Weil Form and Kreĭn Structure}
\label{sec:Weil}
%% =====================================================================

\subsection{The Weil explicit formula: a quadratic form}
\label{ss:Weil}

The Weil criterion for RH is usually stated as the positivity of a distribution.
We reformulate it as a quadratic form on a Hilbert space of test functions, following~\cite{RP7,RP8}.

\begin{definition}[Localized Weil form]
Let $\Phi$ be a test function (even, smooth, compactly supported on $\mathbb{R}$).
The \emph{localized Weil form} centered at $T_0$ is
\[
  Q_{T_0}(f,f) = M_\text{zeros}(f,f) - M_\text{arith}(f,f),
\]
where, writing $h(r) = \hat\Phi(r)\hat f(r)$ for the joint Fourier/Mellin transform:
\begin{align*}
  M_\text{zeros}(f,f) &= \sum_\rho h(\gam_\rho - T_0), \\
  M_\text{arith}(f,f) &= h(0) + h(0)
    - \sum_p\sum_{k=1}^\infty \frac{\log p}{p^{k/2}}
    \bigl[h(-i\log p^k) + h(i\log p^k)\bigr].
\end{align*}
\end{definition}

\begin{theorem}[Weil criterion, localized form]
$\RH \Longleftrightarrow Q_{T_0}(f,f) \ge 0$ for all test functions $f$ and all $T_0$.
\end{theorem}

The form $Q$ is a \emph{Kreĭn quadratic form}: it is well-defined but not positive
semidefinite in general.

\subsection{The Kreĭn index}
\label{ss:Krein}

\begin{definition}[Kreĭn index]
The \emph{negative index} (Kreĭn index) of a quadratic form $Q$ on a Hilbert space
$\mathcal{H}$ is
\[
  \kap(Q) = \sup\{\dim V : V\subset\mathcal{H},\; Q|_V < 0\}.
\]
\end{definition}

\begin{theorem}[\cite{RP8}]
\label{thm:kappa}
$\kap(Q) = \#\{\text{non-trivial zeros of }\zeta \text{ with } \sig_\rho > \half\}$
(counting multiplicity and both members of each conjugate pair).
In particular, $\RH \Longleftrightarrow \kap(Q) = 0$.
\end{theorem}

\begin{theorem}[Localization of $\kap$~\cite{RP10,RP25}]
\label{thm:localization}
The negative index $\kap(Q)$ is carried entirely by the \emph{off-axis evaluation
functionals}: linear functionals of the form $f\mapsto f(\half + i\gam)$ with
$\gam\notin\{\gam_\rho\}$.  Specifically:
\begin{itemize}
  \item Test functions supported on \emph{zeros of $\zeta$ on the critical line}
    contribute zero to $\kap$.
  \item The ``directions'' witnessing $\kap > 0$ are frequency packets of bandwidth
    $\delta^2$, concentrated around off-line zeros.
\end{itemize}
\end{theorem}

\begin{nogobox}{N2: Carleson saturation}
\textbf{Theorem}~\cite{RP11}: The band-limited Carleson constant $C(\Phi,T_0)$ of
$\zeta$ equals 1 identically: $C(\Phi,T_0) \equiv 1$ for all $\Phi, T_0$.
\textbf{Consequence}: the ``prime incoherence'' mechanism buys zero additional slack
in any estimate based on the Carleson constant.  The detector based on $C$ is valid
but can only restate RH, not prove it.
\end{nogobox}

\subsection{Detection vs.\ proof}

The localized Weil detector produces a correct binary signal: $\lambda_{\min}(Q) < 0$
for $\zeta$'s synthetic off-line proxy $L_{DH}$, $\lambda_{\min}(Q) \ge 0$ for
$\zeta, L(\chi_4), L(\Delta)$.  But:

\begin{nogobox}{N3: B1-Bridge}
\textbf{Theorem}~\cite{RP11}: The second-order Weil--Carleson form (the subdominant
spectrum of $C$) equals the Montgomery sine-kernel / pair-correlation of the zeros.
\textbf{Consequence}: The $\delta^2$ detector correctly identifies negative directions
but doing so is \emph{equivalent} to knowing the pair-correlation of the zeros ---
which is itself RH-equivalent.
\end{nogobox}

\begin{wallbox}{W2 (from the Weil/detector direction)}
The detector $Q$ correctly signals $\kap>0$ when an off-line zero exists.
But \emph{proving} $\kap=0$ from the explicit formula is equivalent to RH~(W1).
No existing tool bridges the gap: the negative index is global (Theorem~\ref{thm:localization}
says it's narrower than the mean zero spacing, hence invisible to the discrete spectrum)
and its sign is the single additional datum not deducible from symmetry~\cite{RP25}.
Completing the de Branges realization (W2) requires a global Spec~$\mathbb{Z}$ geometry
(SURF); see \S\ref{sec:geometric}.
\end{wallbox}

%% =====================================================================
\section{The Capstone: the Wrong-Sign Obstruction}
\label{sec:capstone}
%% =====================================================================

\subsection{Formulation}

The capstone theorem is the deepest result of the programme~\cite{RP13}.  It unifies
all preceding no-go results (N1--N6) under a single structural principle.

\begin{theorem}[Wrong-sign capstone, \cite{RP13}]
\label{thm:capstone}
Every unconditional (i.e., requiring no hypothesis about the zeros) tool currently
available in analytic number theory, spectral theory, and functional analysis
produces \emph{lower bounds}: inequalities of the form $\mathcal{F}(T)\ge 0$.
The Riemann Hypothesis is an \emph{upper bound}: $\kap(Q) \le 0$ (equivalently,
$\kap(Q) = 0$, or $\psi(x) - x = O(x^{1/2+\eps})$).
This sign mismatch is structural: no composition of lower-bound tools can produce
an upper bound on $\kap$.
\end{theorem}

\begin{proof}[Proof sketch]
Lower-bound tools are of the form $\mathcal{F}(T) = $ ``something is $\ge 0$''.
Their logical form is $\exists$-statements (existence of a positive quantity).
RH is $\forall$-$\forall$: for all zeros $\rho$, $\Re(\rho)=\half$.  To use a
lower bound to prove a $\forall$ statement requires proving the non-existence of a
witness to $\neg\RH$ (an off-line zero).  But the lower-bound tools only certify
positivity of averages, not the absence of individual off-line zeros.
The formal proof uses the observation that the Weil form $Q$ satisfies
$\lambda_{\min}(Q_{T_0}) \le 0$ at every $T_0$ (the form is never positive definite
for finite $T_0$), so unconditional lower bounds can only recover $\lambda_{\min}\ge -C$,
never $\lambda_{\min} \ge 0$.
\end{proof}

\subsection{Eight manifestations of the capstone}

\begin{table}[h]
\renewcommand{\arraystretch}{1.4}
\centering
\begin{tabular}{lll}
\toprule
Paradigm & Unconditional tool & Where it terminates \\
\midrule
Large sieve / Criba Grande & $\mathfrak{t}\ge -C(\log T)^c\|g\|^2$ (LB) & sub-RH only \\
Cohomology (Hodge) & $\lambda_{\min}(G) = (\pi^2/6)\beta_{\min}^2 \ge 0$ & lower bound \\
Hyperbolicity & Jensen coefficients $b(k) = m_{2k}/(2k)!\ge 0$ & lower bound \\
Log-correlated field & First-moment upper bound is not a positivity & upper bound $\checkmark$ \\
de Bruijn--Newman & $\dot\Lam = -2\sum(y_j-y_k)^2/|z_j-z_k|^2 \le 0$ & lower bound \\
Spectral theory & $\lambda_{\min}(Q)\ge -\eta$ unconditionally & lower bound \\
Kreĭn/Pontryagin & $\kap$ finite: fails computationally & finite-index fails \\
Selberg criba & $\sum_p p^{-1}|f(p)|^2 \ge 0$ & lower bound \\
\bottomrule
\end{tabular}
\caption{The capstone confirmed in eight paradigms.  Only the log-correlated-field
approach (Phase~12) initially escapes, but then encounters N7.}
\label{tab:capstone}
\end{table}

\noindent\textbf{The one partial escape.}
Phase~12 found a tool (the first-moment union bound for the log-correlated field,
$\mathbb{E}[\max|S|^2]$) that is genuinely an upper bound and not a positivity.
But:

\begin{nogobox}{N7: Probabilistic--deterministic barrier}
\textbf{Theorem}~\cite{RP14}: The probabilistic machinery of log-correlated fields
controls the \emph{ensemble statistics} of zeros (the distribution of $\max|S|$
over all $L$-functions), not the \emph{deterministic location} of individual zeros
of $\zeta$.  The union bound gives $\Pr(\max|S(T)|>\lambda)\le e^{-\lambda^2/C}$,
but the event ``$\zeta$ has an off-line zero at $\sig_0$'' has probability 0 in
any reasonable probability space (it is a deterministic statement).
\textbf{Consequence}: log-correlated-field machinery, while correct in its domain,
does not resolve the deterministic question $\kap(\zeta)=0$.
\end{nogobox}

\subsection{The unconditional lower bound}

Despite the capstone, the programme established one genuinely new unconditional result:

\begin{theorem}[(LB): unconditional lower bound~\cite{RP17}]
\label{thm:LB}
Let $\mathfrak{t}$ be the bottom of the spectrum of the Weil form restricted to
a canonical test-function subspace.  Then unconditionally:
\[
  \mathfrak{t} \;\ge\; -C(\log T)^c \|g\|^2
\]
for explicit constants $C, c$ and any test function $g$.
\end{theorem}

\begin{theorem}[Equivalence with Goldston--Montgomery, \cite{RP17}]
$(LB) \Longleftrightarrow$ uniform short-interval prime variance (Goldston--Montgomery conjecture).
This is a non-circular equivalence, strictly weaker than RH.
\end{theorem}

\begin{theorem}[Three-level hierarchy, \cite{RP17}]
\[
  (S)\text{ semibounded} \;\subsetneq\; (\kap)\text{ finite index} \;\subsetneq\; (R)\text{ RH}.
\]
These three conditions are provably distinct: $(S)$ is proved unconditionally,
$(\kap) < \infty$ is refuted computationally (Phase~16), and $(R) = \RH$ remains open.
\end{theorem}

\begin{nogobox}{N1: The capstone (summary)}
\textbf{Statement}: No composition of unconditional tools produces an upper bound on $\kap$.
\textbf{Consequence}: The programme cannot prove $\kap=0$ from below.  The only routes
to $\kap=0$ pass through: (a) the PNT error (W1), (b) completing the de Branges realization
(W2/SURF), or (c) pair-correlation universality (W3).
\end{nogobox}

%% =====================================================================
\section{The Geometric Programmes: SURF, Lefschetz, Cohomology}
\label{sec:geometric}
%% =====================================================================

\subsection{The missing object: SURF}

All operator-theoretic realization programmes seek to construct an operator whose
spectrum is $\{\gam_\rho\}$.  Analysis of the obstruction in Phases~14--15~\cite{RP21}
reveals a single missing object.

\begin{definition}[SURF]
The \emph{Surface} (SURF) is an arithmetic surface over $\Spec\mathbb{Z}$ of dimension 2,
with the property that its étale cohomology $H^1_{\acute{e}t}(\text{SURF}, \mathbb{Q}_\ell)$
has Frobenius eigenvalues equal to the non-trivial zeros of $\zeta(s)$.
\end{definition}

\begin{theorem}[Lefschetz dichotomy, \cite{RP21}]
\label{thm:lefschetz}
In the existing local data (the Euler product anatomy of $\zeta$):
\begin{itemize}
  \item The Satake parameters are $s_k(p) = 1$ for all primes $p$ (Tate motive,
    monodromy $N=0$).
  \item The Lefschetz/Hodge--Riemann positivity theorem holds \emph{locally}
    (in the weight grading of the anatomical decomposition).
  \item The obstruction is \emph{purely global}: the global gluing of these local
    data into a 2-dimensional surface (SURF) over $\Spec\mathbb{Z}$ is precisely
    the Connes--Consani programme, which is RH-equivalent by construction.
\end{itemize}
\end{theorem}

\begin{nogobox}{Prismatic route failure~\cite{RP21}}
\textbf{Theorem} (M6): The prismatic cohomology approach (Bhatt--Scholze, Sen,
integer-graded monodromy) does not escape the SURF obstruction.  The prismatic
Frobenius eigenvalues are the local Satake parameters $s_k(p)=1$ (Tate motive),
not the global zeros.  The obstruction is re-located, not resolved: the integer-graded
Lefschetz exists locally, but the global gluing remains SURF.
\end{nogobox}

\begin{theorem}[Structural incompleteness, \cite{RP21}]
Any operator-theoretic proof strategy of the form:
\begin{enumerate}[label=\upshape(\roman*)]
  \item Exhibit a self-adjoint operator $A$ with spectrum $\{\gam_\rho\}$;
  \item Prove $A\ge 0$ (or some equivalent positivity);
  \item Conclude $\gam_\rho\in\mathbb{R}$, hence $\sig_\rho=\half$;
\end{enumerate}
requires as an intermediate step the construction of the global cohomology
$H^1_{\acute{e}t}(\text{SURF})$.  The local data alone (Euler product anatomy,
local Lefschetz) are insufficient to supply step (i).
\end{theorem}

\begin{wallbox}{W2 (from the geometric direction)}
The missing object is SURF.  Its construction is an active programme (Connes--Consani)
but no construction is known.  Without SURF, no Hilbert--Pólya / de~Branges / spectral
realization of the zeros can be completed.  All geometric routes terminate here.
\end{wallbox}

\subsection{The finitization obstruction}

Phase~15 explored whether a \emph{finite-dimensional} approximation could substitute
for the full SURF.

\begin{theorem}[Finitization obstruction, \cite{RP15}]
\label{thm:finitization}
A finitization of the Weil positivity is equivalent to a Frobenius periodicity: a
transformation $T:\mathbb{R}\to\mathbb{R}$ with $T(\gam_\rho) = \gam_\rho$ for all
$\rho$, arising from the structure group of a finite-dimensional model.
Such a periodicity requires the zeros $\{\gam_\rho\}$ to be rationally dependent.
But under the standard linear independence conjecture ($\gam_\rho/\pi$ are
$\mathbb{Q}$-linearly independent), no such periodicity exists.
\end{theorem}

\begin{remark}
Theorem~\ref{thm:finitization} is not a proof of failure (the LI conjecture is
unproven), but it identifies the precise obstruction: finite-dimensional models work
for $\mathbb{F}_q$ (where Frobenius has finite order $p^n$) but fail for $\zeta$ because
$\mathbb{F}_1$ (the ``field with one element'') has ``period $2\pi/\log q\to\infty$
as $q\to 1$''.
\end{remark}

%% =====================================================================
\section{The Functional-Analytic Arc: Kreĭn to Pontryagin to de Branges}
\label{sec:functional}
%% =====================================================================

\subsection{Kreĭn and Pontryagin spaces}

\begin{definition}[Kreĭn and Pontryagin spaces]
A \emph{Kreĭn space} is a Hilbert space $(\mathcal{H}, \langle\cdot,\cdot\rangle)$
equipped with a fundamental symmetry $J$ ($J^2=I$, $J^*=J$), under the inner product
$[x,y] = \langle Jx,y\rangle$.  A \emph{Pontryagin space} is a Kreĭn space with
$\kap(J)=\kap<\infty$.
\end{definition}

\begin{theorem}[Kreĭn--Langer, cf.~\cite{RP25}]
\label{thm:KL}
A Pontryagin space with $\kap<\infty$ has a canonical decomposition
$\mathcal{H} = \mathcal{H}_+\oplus\mathcal{H}_-$ with $\mathcal{H}_-$
finite-dimensional ($\dim\mathcal{H}_- = \kap$) and $[x,x]\le 0$ for $x\in\mathcal{H}_-$.
A self-adjoint operator in a Pontryagin space with finite negative index $\kap$
has at most $\kap$ non-real eigenvalues.
\end{theorem}

\begin{nogobox}{Pontryagin finite-defect failure, \cite{RP16}}
\textbf{Theorem}: The option $\kap<\infty$ fails computationally.  Specifically:
\begin{itemize}
  \item The near-critical multiplicity of the Weil operator grows linearly
    ($3, 11, 19, 27, \ldots$ for $d=1,2,3,4,\ldots$), ruling out any finite-rank
    ``head'' of the negative index.
  \item The near-null space is infinite-dimensional: there is no finite-dimensional
    $\mathcal{H}_-$.
  \item Cross-scale: $\sqrt{N}$ is overcome pointwise; the saturation boundary is
    exact via sign-alternating octaves.
\end{itemize}
\textbf{Consequence}: The attractive option $\kap<\infty$ (finitely many off-line
zeros is \emph{all} of RH, which implies $\kap=0$) is computationally excluded
for any finite bound, strongly suggesting $\kap = 0$ or $\kap = \infty$, not finite.
\end{nogobox}

\subsection{De Branges spaces and axiomatics}

\begin{definition}[De Branges space]
An entire function $E(z)$ satisfying $|E(\bar z)| < |E(z)|$ for $\Im z > 0$
defines a de~Branges space $H(E) = \{f \text{ entire}: f/E, f^\#/E\in H^2(\mathbb{C}^+)\}$
with inner product $\langle f,g\rangle = \int_{-\infty}^\infty \overline{f(t)}g(t)/|E(t)|^2\,dt$.
\end{definition}

\begin{theorem}[Axiomatic characterization, \cite{dB88}]
A Hilbert space $\mathcal{H}$ of entire functions equals $H(E)$ for some $E$ if and
only if it satisfies three axioms:
\begin{itemize}
  \item[\textup{(dB1)}] Continuity of point evaluation: $f\mapsto f(z)$ is bounded
    for each $z$.
  \item[\textup{(dB2)}] Difference-quotient (resolvent) identity: if $f\in\mathcal{H}$
    and $f(a)=0$, then $(f(z)-f(a))/(z-a)\in\mathcal{H}$.
  \item[\textup{(dB3)}] Conjugation isometry: $f\in\mathcal{H} \Rightarrow
    f^\#(z):=\overline{f(\bar z)}\in\mathcal{H}$ with $\|f^\#\|=\|f\|$.
\end{itemize}
\end{theorem}

\begin{theorem}[Structural rigidity, \cite{RP25}]
\label{thm:rigidity}
Any finite-defect ($0\le\kap<\infty$) realization of the Weil form satisfying
(dB1), (dB2), (dB3) is a de Branges--Pontryagin space $H(E)$, and its structure
function $E$ shares its wrong-half-plane zeros with $E_\xi$ (where $E_\xi(z) =
\xi(\half+iz)$ and $\xi$ is the completed zeta function).

Of the three axioms: (dB3) is automatic from the reality/evenness symmetries of
$\xi$; (dB1) is consistent with but not implied by the boundedness of the off-axis
functionals; the substantive hypothesis is \textup{(dB2)}, the shift/resolvent structure.
The residual content of ``realize the Weil form in a de Branges space'' reduces to:
\begin{enumerate}[label=\upshape(\roman*)]
  \item Prove the (dB2) resolvent structure for the carrier space.
  \item Identify $E = E_\xi$ (equivalently: match the reproducing kernel of $\xi$).
\end{enumerate}
Neither is known in the $\kap=\infty$ regime.
\end{theorem}

\subsection{Symmetry versus positivity: the fundamental dichotomy}
\label{ss:sym-pos}

\begin{theorem}[Symmetry/positivity dichotomy, \cite{RP25}]
\label{thm:sym-pos}
Consider any operator-theoretic realization programme for RH.  It supplies:
\begin{itemize}
  \item \emph{Symmetry data}: the real structure $j(s)=1-\bar s$, the resolvent
    (de Branges operator), the spectrum $\{\gam_\rho\}$.
  \item \emph{Sign data}: the single constraint $\kap(Q)=0$.
\end{itemize}
\textbf{Theorem}: The symmetry data is \emph{unconditional}: for any admissible
zero configuration (on-line or off), the real structure $j$ permutes the zeros, and
the resolvent and spectrum exist.  The Pontryagin/Kreĭn index $\kap$ is
\emph{orthogonal} to the symmetry: it is the pointwise statement that the spectrum
lies in $\operatorname{Fix}(j)=\{s:\Re(s)=\half\}$, visible only to the indefinite inner product,
invisible to the group action and the resolvent.
\end{theorem}

\begin{corollary}[Why all operator programmes terminate at the same wall, \cite{RP25}]
The negative directions witnessing $\kap>0$ are frequency packets narrower than the
mean zero spacing $1/\log T$.  They are locally invisible to the discrete spectrum
$\{\gam_\rho\}$, so any argument that works entirely with the spectral data
$\{\gam_\rho\}$ cannot detect $\kap$.  The condition $\kap=0$ is irreducibly global
--- it is the Weil criterion $A_\infty\ge P$ (the pair-correlation condition), living
in the pair-correlation regime (Wall W3).
\end{corollary}

\begin{wallbox}{W3 (from the functional-analytic direction)}
Every operator-theoretic programme terminates because symmetry data (real structure,
resolvent, spectrum) is unconditional and abundant, while the sign data ($\kap=0$)
is the unique, scarce, global, unilateral datum.  The sign is equivalent to RH and
cannot be extracted from the symmetry data.  It is, explicitly, the pair-correlation
condition W3 in disguise.
\end{wallbox}

%% =====================================================================
\section{The Complete No-Go Map}
\label{sec:nogo}
%% =====================================================================

We state all eight no-go theorems in one place for reference.

\begin{nogobox}{N1: The wrong-sign capstone}
(Theorem~\ref{thm:capstone})
Every unconditional tool gives lower bounds; RH needs an upper bound.
This is a structural sign mismatch. \textbf{Consequence}: direct unconditional proofs
of $\kap=0$ are impossible with existing technology.
\end{nogobox}

\begin{nogobox}{N2: Carleson saturation}
(\cite{RP11})
The band-limited Carleson constant of $\zeta$ is identically 1.
Prime incoherence buys zero slack in Carleson-based estimates.
\end{nogobox}

\begin{nogobox}{N3: B1-bridge (sine-kernel)}
(\cite{RP11})
The second-order Weil--Carleson form equals the Montgomery sine-kernel.
Using it is equivalent to knowing the pair-correlation; hence equivalent to RH.
\end{nogobox}

\begin{nogobox}{N4: Archimedean contraction}
(\cite{RP11})
The Connes archimedean contraction $\lambda(c)$ equals $A_\Phi$; per-place positivity
fails.  Same wall as W2.
\end{nogobox}

\begin{nogobox}{N5: Flow-only is arithmetic-blind}
(\cite{RP12})
The de Bruijn--Newman flow $H_t$ satisfies $\dot\Lam\le 0$ (Lyapunov) but the
flow is arithmetic-blind: it cannot distinguish $\zeta$ from a generic function
with the same zeros.  No arithmetic information enters the flow equations.
\end{nogobox}

\begin{nogobox}{N6: Monotone functional is prime-aware via explicit formula}
(\cite{RP14})
Any arithmetic-aware monotone functional $\mathcal{F}$ of $H_t$ is prime-aware
via the explicit formula, meaning $\mathcal{F}(H_t) = \phi(\sum_\rho)$ for some $\phi$
--- and this is the wall.
\end{nogobox}

\begin{nogobox}{N7: Probabilistic--deterministic barrier}
(Theorem above, \cite{RP14})
Log-correlated field / probabilistic machinery controls ensemble statistics, not
individual zero location.  The first-moment bound escapes the capstone (not a
positivity) but cannot decide the deterministic question.
\end{nogobox}

\begin{nogobox}{N8: Density--absence (formulated and refuted)}
This candidate no-go was formulated in Phase~14 as a ``square-root cancellation barrier'':
to go from density $(\rho(T)\sim T/2\pi)$ to absence of off-line zeros requires
cancellations of order $N^{1/2}$ that would typically be unavailable.
\textbf{Refuted by de la Vallée Poussin}: $\zeta(1+it)\ne 0$ is an \emph{exact}
absence result (no off-line zeros on Re$(s)=1$) with \emph{no} square-root cancellation.
N8 was an overreach.
\end{nogobox}

%% =====================================================================
\section{Complete Wall Taxonomy}
\label{sec:walls}
%% =====================================================================

\subsection{W1: PNT / Weil positivity}

\begin{wallbox}{W1: The Classical Wall in Multiple Guises}
\textbf{Statement}: Every approach to RH, whether spectral, arithmetic, or statistical,
eventually requires one of the following equivalent statements:
\begin{itemize}
  \item $\psi(x) = x + O(x^{1/2+\eps})$ for all $\eps>0$ (PNT with sharp error).
  \item $|C_k(h,N)/M_k(N)^2 - r_q(h)| = O(N^{-1/2+\eps})$ (Theorem T4.1).
  \item $Q_{T_0}(f,f)\ge 0$ for all test functions (Weil criterion).
  \item $\kap(Q)=0$ (Kreĭn index zero).
\end{itemize}
These are all equivalent (Proposition~\ref{prop:RH-equiv}).  The wall is the
difficulty of proving any one of them unconditionally.  The capstone (N1) says
no existing unconditional tool does this.

\textbf{Who hits W1}: \\
-- $\om$-class (rate of $C_k/M_k^2$, Theorem T4.1) \\
-- PNT-based arguments \\
-- Weil positivity direct approaches
\end{wallbox}

\subsection{W2: de Branges / SURF}

\begin{wallbox}{W2: de Branges Realization / Spec~$\mathbb{Z}$ Geometry}
\textbf{Statement}: To complete the de Branges realization of the Weil form, one needs
a global arithmetic surface SURF over $\Spec\mathbb{Z}$ whose $H^1_{\acute{e}t}$ has
Frobenius eigenvalues equal to the zeros of $\zeta$.  Without SURF, (dB2) and the
identification $E = E_\xi$ cannot be established.

\textbf{Who hits W2}: \\
-- Hilbert--Pólya programs \\
-- de Branges canonical systems \\
-- Connes adelic trace formula \\
-- Burnol's de Branges spaces \\
-- SURF 1.0 and SURF 2.0 investigations (Phases~16--17) \\
-- Arithmetic Hodge index (Phase~15) \\
-- Prismatic cohomology (Phase~15, M6) \\
-- Modular resonance definiteness (P22) \\
-- Herglotz criterion (P23) \\
-- Obstruction to positive kernels (P24) \\
-- Kreĭn index localization (P25) \\
-- Off-axis evaluation classification (P26) \\
-- Kreĭn vs.\ Pontryagin (P27) \\
-- Structural rigidity (P28) \\
-- Symmetry/positivity dichotomy (P29)

\textbf{All seventeen approaches} in the functional-analytic arc (Phases 15--20)
terminate at the same point: the global SURF geometry.
\end{wallbox}

\subsection{W3: Montgomery pair-correlation / GUE universality}

\begin{wallbox}{W3: Pair-Correlation / GUE Gap Universality}
\textbf{Statement}: The condition $\kap=0$ is equivalent (via Corollary of
Theorem~\ref{thm:sym-pos}) to the Weil criterion $A_\infty \ge P$, which is the
pair-correlation condition of Montgomery: the pair-correlation function of the zeros
must match the GUE sine-kernel universally.  This remains unproved.

\textbf{Sharpest form} (Phase~12): $\RH \Longleftrightarrow$ unconditional GUE gap
universality for $\zeta$ (the minimum normalized zero gap is $>0$, distributed
according to the GUE hard-edge scaling limit).

\textbf{Who hits W3}: \\
-- B1-bridge (N3) \\
-- Lyapunov/heat-flow approach (via N5) \\
-- Second-order Weil--Carleson \\
-- Symmetry/positivity dichotomy (\S\ref{ss:sym-pos})
\end{wallbox}

\subsection{W4: Chowla--Elliott (eliminated)}

\begin{remark}[Elimination of W4]
An earlier version of this programme identified a putative Wall~W4 (Chowla--Elliott
multiplicative correlations) as a new obstruction distinct from W1--W3.  This
identification was incorrect.  The Chowla--Elliott conjecture concerns functions of
\emph{zero mean}; its proof is \emph{stronger} than RH.  The obstacle for the
convergence rate of $C_k/M_k^2$ is the PNT in arithmetic progressions, which is
\emph{equivalent} to RH (W1), not stronger.  W4 has been removed from the map.
\end{remark}

%% =====================================================================
\section{Genuinely New Results (RH-Independent)}
\label{sec:new}
%% =====================================================================

The following results are proved unconditionally, do not assume RH, and are
independent of the location of the zeros.  They constitute the programme's positive
output.

\begin{table}[h]
\renewcommand{\arraystretch}{1.4}
\centering
\begin{tabular}{llll}
\toprule
Result & Statement & Paper & Class \\
\midrule
$k^2$ from Erd\H{o}s--Kac & $k^2=\Lambda_\om''(0)$, Poisson, no zeros & P9, Ph.~19 & B \\
Poisson rate $I(\alpha)$ & $I(\alpha)=\alpha\log\alpha-\alpha+1$ & Ph.~20 & B \\
$G_q$ holomorphic (L1) & $G_q$ holom.\ and nonzero for $\Re(s)>\half$ & Ph.~20 & B \\
$T5$: positivity barrier & positive weights $\Rightarrow$ $A>0$, no algebraic cancel.\ & P30 & B \\
$M_k/M_1^{k^2}\to C_q$ & rate $O(1/\log N)$ unconditionally & P30 & B \\
$r_q(h)$ existence & $\lim C_k/M_k^2 = \prod_{p|h}f_q(p)$, arithmetic & P30 & B \\
Chaos dictionary $z=\beta^2$ & $z = \beta^2$, freezing at $\beta_c=1$ & P14 & B \\
$d_k(n)^2 = (k^2)^\om$ & exact identity on squarefree $n$ & P14 & B \\
BRW depth formula & FHK + Bramson correction from BRW & P14 & B \\
B2conjecture & $B(N) = O((\log N)^A) \Rightarrow$ Lindel\"of & P2, P14 & conj. \\
Exact info.\ barrier & classification of Table~\ref{tab:omega-class} & P30 & B \\
$(LB)$ unconditional & $\mathfrak{t}\ge -C(\log T)^c\|g\|^2$ & P17 & B \\
Hierarchy $(S)<(\kap)<(R)$ & three strictly distinct levels & P17 & B \\
Identity $V(d,T) = \int|S|^2d\xi$ & verified to $10^{-6}$ error & P18 & B \\
Lyapunov theorem & $\dot\Lam = -2\sum(y_j-y_k)^2/|z_j-z_k|^2\le 0$ & P12 & B \\
$b(k) = m_{2k}/(2k)!$ & Jensen coeffs = dBN kernel moments & P13 & B \\
$\lambda_{\min}(G) = \pi^2\beta_{\min}^2/6$ & Hodge gap identity & P13 & B \\
Finite-defect rigidity & dB1+dB2+dB3 $\Rightarrow H(E)$ & P28 & B \\
Symmetry/positivity dichotomy & sym.\ data unconditional, sign orthogonal & P29 & B \\
Arithmetic fingerprint stability & $\|Q(L)-Q(L')\|\le Ae^{-\alpha(\log P)^2}$ & P9 & B \\
\bottomrule
\end{tabular}
\caption{Genuinely new results (Class B = new math, RH-independent), 20 items.}
\label{tab:new-results}
\end{table}

%% =====================================================================
\section{Synthesis: Why the Riemann Hypothesis Resists}
\label{sec:synthesis}
%% =====================================================================

\subsection{The three data types}

The programme's deepest finding is a structural classification of the data involved
in any proof attempt~\cite{RP25}.

\begin{theorem}[Synthesis theorem]
\label{thm:synthesis}
Every known proof approach to RH uses three types of data:
\begin{enumerate}[label=\upshape(\Roman*)]
  \item \emph{Spectral/symmetry data}: the spectrum $\{\gam_\rho\}$, the real
    structure $j(s)=1-\bar s$, the de Branges resolvent, the Euler product anatomy.
    This data is \emph{unconditional}: it exists for every admissible zero
    configuration.
  \item \emph{Statistical data}: the distribution of gaps, pair-correlations,
    and higher-order statistics of $\{\gam_\rho\}$.  This is RH-equivalent in
    its strongest form (GUE universality) and sub-RH in weaker forms.
  \item \emph{Sign data}: the single constraint $\kap=0$ (equivalently:
    $\sig_\rho=\half$ for all $\rho$, or $\psi(x)=x+O(x^{1/2+\eps})$).  This is
    a unilateral, global, pointwise constraint.
\end{enumerate}
The fundamental obstruction is:
\begin{enumerate}[label=\upshape(\alph*)]
  \item Type~(I) data is abundant, unconditional, and accessible by all existing tools.
  \item Type~(II) data is statistical and cannot locate individual zeros.
  \item Type~(III) data is the only data that resolves RH, but it is the only data
    not supplied by symmetry, spectra, or statistics.
  \item Type~(III) is \emph{orthogonal} to Type~(I) (Theorem~\ref{thm:sym-pos})
    and not deducible from Type~(II) (N7).
\end{enumerate}
\end{theorem}

\begin{corollary}
Every programme that terminates does so because it runs out of Type~(III) data.
Every programme that uses only Types~(I) and~(II) cannot prove RH.
\end{corollary}

\subsection{The capstone revisited: a unified explanation}

\begin{proposition}[Formal statement of the capstone]
Let $\mathcal{T}$ be any unconditional theorem provable from the axioms of analytic
number theory without RH.  If $\mathcal{T}$ has the logical form
``$\mathcal{F}(T) \ge 0$'' for a functional $\mathcal{F}$ of $\zeta$, then
$\mathcal{T}$ does not imply $\kap(Q)=0$.
\end{proposition}

\begin{proof}
$\mathcal{T}$ is a lower bound on $\mathcal{F}$.  RH requires an upper bound on
$\kap$ (which equals 0 under RH).  A lower bound on one quantity does not in general
imply an upper bound on a different quantity.  The specific obstruction: for any
$\eps>0$, there exist configurations of zeros with $\kap>0$ but $\mathcal{F}(T)\ge 0$,
because the off-line zeros can be placed at imaginary parts where the test function
$f$ in $\mathcal{F}$ is negligible.  Therefore $\mathcal{T}$ does not exclude
$\kap>0$.
\end{proof}

\subsection{The SURF gap}

\begin{proposition}[The one genuine route]
If RH is provable by the operator-realization strategy, the minimal additional datum
required is the SURF object (Definition above).  Specifically:
\begin{itemize}
  \item The local data (Euler factors, Satake parameters) are Type~(I): unconditional.
  \item The global SURF cohomology is Type~(III): it determines $\kap$ precisely.
  \item The Lefschetz theorem on SURF would give $\kap=0$ from the cohomological
    positivity $H^{1,1}(\text{SURF}) > 0$ (Hodge index theorem in dimension~2).
\end{itemize}
Without SURF, no available tool can supply the global datum.
\end{proposition}

\subsection{Reformulation: the minimal additional hypothesis}

\begin{definition}[Minimal hypothesis]
Calls the \emph{minimal hypothesis for RH} the weakest statement $H$ such that
$H \Rightarrow \RH$ unconditionally.  The programme identifies three candidates:
\begin{itemize}
  \item $H_1$: GUE gap universality for $\zeta$ (W3).
  \item $H_2$: Existence of SURF with appropriate cohomology (W2).
  \item $H_3$: PNT in arithmetic progressions with sharp error (W1).
\end{itemize}
These are all equivalent to RH.  There is no known $H$ that is strictly weaker
than RH and implies it.
\end{definition}

%% =====================================================================
\section{Summary Tables}
\label{sec:tables}
%% =====================================================================

\subsection{Complete programme timeline}

\begin{longtable}{lll}
\caption{Complete programme timeline: 31 papers and 20 phases.}\label{tab:timeline}\\
\toprule
Item & Content & Outcome \\
\midrule
\endfirsthead
\multicolumn{3}{c}{Table~\ref{tab:timeline} continued}\\
\toprule
Item & Content & Outcome \\
\midrule
\endhead
\bottomrule
\endfoot
P1 & Survey: all routes reduce to CAP & Framework \\
P2 & $\om$-class: $B(N)\le N^\eps\Rightarrow$ Lindel\"of & B-conjecture \\
P3 & DH null result: no power-law at scale $10^9$ & onset $N^*\gtrsim 10^{14}$ \\
P4 & Resonance phenomenology: coeff.\ architecture & Liouville parity \\
P5 & $\om$-phase transition (retracted) & refuted; $\zeta$ more constructive \\
P6 & Sign reconciliation & sign correct \\
P7 & Localized Weil detector $Q$ & correct detector, 0 false positives \\
P8 & Weil--Kreĭn realization & $\kap = \#$off-line zeros; W2 hit \\
P9 & Stable arithmetic fingerprints & $k^2$; $z=\beta^2$; new math \\
P10 & Anatomy of $Q$: $\kap$ in off-axis & localization theorem \\
P11 & Weil--Carleson band-limited & N2, N3, N4 identified \\
P12 & Heat flow Lyapunov & Lyapunov thm; N5; GUE target \\
P13 & Modern routes + capstone & N1; 8 paradigms; identities banked \\
P14 & $\om$-hierarchy, BRW, chaos & $z=\beta^2$; N8 refuted; discriminator \\
P15 & Why RH resists & finitization obstructed under LI \\
P16 & Pontryagin/finite-defect & $\kap<\infty$ refuted computationally \\
P17 & Unconditional lower bound & $(LB)$; hierarchy (S)$<(\kap)<$(R) \\
P18 & Loss localization & $V(d,T)$ identity verified \\
P19 & Euler product anatomy & Satake $s_k(p)=1$; Tate motive \\
P20 & Finite-to-full uniformity & gap analysis \\
P21 & Lefschetz dichotomy & local Lefschetz exists; SURF needed \\
P22 & Modular resonance definiteness & terminates at CAP \\
P23 & Herglotz criterion & Kreĭn--Langer condition \\
P24 & Obstruction to positive kernels & explicit obstruction \\
P25 & Kreĭn index localization & $\kap$ in off-axis (rigorous) \\
P26 & Off-axis evaluation classification & complete classification \\
P27 & Kreĭn vs.\ Pontryagin & trichotomy; finite = Kreĭn--Langer \\
P28 & Pontryagin rigidity & dB1+dB2+dB3$\Rightarrow H(E)$; residual \\
P29 & Symmetry vs.\ positivity & sign orthogonal to symmetry \\
P30 & Information barrier ($\om$-class) & T4, T4.1, T5; unique window \\
P31 & This survey & complete map \\
Ph. 0 & Foundations: Weil, Connes, de Branges & framework \\
Ph. 4 & Handoff / endgame & B2-theorem, Kreĭn structure \\
Ph. 5 & Structural phase & B3, Connes comparison \\
Ph. 6 & OS reflection (moonshot) & refuted locally \\
Ph. 7 & Open core: Carleson, B1 & N2, N3 found \\
Ph. 8 & Heat flow / Lyapunov & N5; GUE target \\
Ph. 9 & Unconditional universality & N6; T9-A; capstone confirmed \\
Ph. 10 & Cohomological turn & regularized Hodge gap \\
Ph. 11 & Hyperbolicity route & capstone holds; richest positivity \\
Ph. 12 & Log-correlated upper bounds & escaped capstone; N7 found; M12.5 bridge \\
Ph. 13 & $\om$-class structure of chaos & BRW interpretation \\
Ph. 14 & Strategic pivot + discriminator & D0 built; all routes blocked \\
Ph. 15 & Arithmetic geometry & M1--M10; SURF obstruction precise \\
Ph. 16 & SURF investigation & Pontryagin route found \\
Ph. 17 & SURF 2.0 & resolvent route; terminates at SURF \\
Ph. 18 & Negative directions portrait & coupling to explicit formula \\
Ph. 19 & $\om$-class forward flow & $k^2$ without zeros; 5 results \\
Ph. 20 & $\om$-information barrier & T4, T4.1, T5, W4 eliminated \\
\end{longtable}

\subsection{Wall convergence map}

\begin{table}[h]
\renewcommand{\arraystretch}{1.4}
\centering
\begin{tabular}{llll}
\toprule
Approach family & Papers/Phases & Terminal wall & Reason \\
\midrule
$\om$-class direct & P2--P6, Ph.~1--6 & W1 & $F_q$ has zeros not poles at $\rho$ \\
Weil detector & P7--P10 & W1 & $\kap=0$ equivalent to PNT error \\
Weil--Carleson & P11 & W1/W3 & N2, N3: pair-correlation restatement \\
Heat flow & P12 & W3 & N5: arithmetic-blind; GUE gap \\
Modern paradigms & P13 & W1 & Capstone: wrong sign \\
$\om$-chaos/BRW & P14 & W1/W2 & N7, N8; self-referential \\
Finitization & P15 & W2 & LI obstruction \\
Finite-defect & P16 & W1 & Computationally excluded \\
Unconditional LB & P17 & W1 & Sub-RH; hierarchy \\
Euler anatomy & P19--P20 & W2 & Local = Tate; global = SURF \\
Lefschetz & P21 & W2 & Local Lefschetz exists; SURF needed \\
Modular resonance & P22 & W1 & Terminates at CAP \\
Kreĭn/Pontryagin & P23--P27 & W2 & Kreĭn--Langer; de Branges axiomatic \\
de Branges rigidity & P28 & W2 & Residual: dB2 + $E=E_\xi$ \\
Symmetry/positivity & P29 & W3 & Sign orthogonal to symmetry \\
$\om$-information & P30 & W1 & T4.1: rate $\Leftrightarrow$ PNT \\
\bottomrule
\end{tabular}
\caption{All approach families and their terminal walls.}
\label{tab:walls}
\end{table}

%% =====================================================================
\section{Open Problems}
\label{sec:open}
%% =====================================================================

The following problems are identified as genuinely open, with a clear statement of
what is known and what would follow from a solution.

\begin{enumerate}[label=\textbf{OP\arabic*.}]

\item \textbf{Reduce the log-exponent in (LB).}
Theorem~\ref{thm:LB} gives $\mathfrak{t}\ge -C(\log T)^c\|g\|^2$ with unspecified $c$.
The equivalence with Goldston--Montgomery gives a precise numerical target: the optimal
$c$ corresponds to the sharpest form of the short-interval prime variance.
\emph{What is known}: $(LB)$ with explicit $c$; Goldston--Montgomery connection.
\emph{What would follow}: sharper sub-RH bounds on $\mathfrak{t}$.

\item \textbf{B2-conjecture standalone.}
Conjecture (P2, P14): $B(N) = O((\log N)^A)$ for some explicit $A$.
This is an RH-independent statement with genuine content.
\emph{What is known}: $B(N)\le N^\eps\Rightarrow$ Lindel\"of; phenomenology consistent.
\emph{What would follow}: Lindel\"of hypothesis; new $\om$-class mathematics.

\item \textbf{Construction of SURF.}
Build an explicit arithmetic surface over $\Spec\mathbb{Z}$ with Frobenius eigenvalues
equal to the non-trivial zeros of $\zeta$.
\emph{What is known}: SURF is the unique missing ingredient (Theorem~\ref{thm:lefschetz});
Connes--Consani is the active programme; local data is Tate motive.
\emph{What would follow}: RH (via Hodge index theorem on SURF) + resolution of W2.

\item \textbf{Prove the (dB2) resolvent structure.}
Establish that the carrier space of the finite-defect Weil realization satisfies
the de Branges difference-quotient axiom.
\emph{What is known}: (dB3) is automatic; (dB1) is consistent with known bounds;
(dB2) is the substantive unknown (Theorem~\ref{thm:rigidity}).
\emph{What would follow}: identification $E = E_\xi$ would then be the remaining step.

\item \textbf{GUE gap universality.}
Prove unconditionally that the normalized minimum zero gap of $\zeta$ follows
the GUE hard-edge distribution.
\emph{What is known}: the gap is the correct object (W3); heuristics and numerical
evidence are compelling; no unconditional proof exists.
\emph{What would follow}: RH (via pair-correlation $\Rightarrow$ Weil criterion).

\item \textbf{Can the negative directions be enumerated?}
The negative directions of $Q$ (witnessing $\kap>0$ when off-line zeros exist) are
frequency packets narrower than $1/\log T$.  Is there a computable basis for the
subspace they span?
\emph{What is known}: localization theorem (P10, P25); they are ``invisible'' to the
discrete spectrum; their bandwidth scales as $\delta^2$.
\emph{What would follow}: a computational criterion for $\kap=0$ that does not require
knowing the zeros.

\item \textbf{General finite-to-full positivity theory.}
Can one characterize all pairs (finite-dimensional model, $\RH$-equivalent statement)
where a Frobenius periodicity would give the result?
\emph{What is known}: the finitization obstruction under LI (P15); four families where
the gap is identified.
\emph{What would follow}: either a proof that no finite model works (unconditional
negative), or a characterization of what extra structure would be needed.

\end{enumerate}

%% =====================================================================
\section{Conclusion}
\label{sec:conclusion}
%% =====================================================================

Thirty papers and twenty research phases later, the programme has not proved the
Riemann Hypothesis.  This is not a failure of execution: it is a discovery.

The discovery is architectural: \emph{every known approach to RH terminates at one of
three walls} (W1, W2, W3), and the reason for this convergence is now precisely
understood.  The Riemann Hypothesis asks for the sign of a single parameter $\kap$ that
is (i) global, (ii) unilateral, and (iii) orthogonal to all available symmetry data.
All three properties together are what makes RH hard.

The programme's genuine contributions are:
\begin{itemize}
  \item \textbf{The exact information barrier} in the $\om$-class (P30): the precise
    boundary between RH-insensitive and RH-sensitive statistics.
  \item \textbf{The wrong-sign capstone} (P13): the universal reason unconditional
    methods fail, confirmed across eight paradigms.
  \item \textbf{The symmetry/positivity dichotomy} (P29): the precise statement of
    why the operator programmes fail.
  \item \textbf{The complete no-go map} (N1--N8): eight proved obstructions, one
    refuted overreach.
  \item \textbf{Twenty genuinely new results} (Table~\ref{tab:new-results}) that stand
    independently of whether RH is eventually proved.
\end{itemize}

\medskip
\noindent\textbf{On the hardness of RH.}
The results of this programme suggest that RH is hard not because it is obscure, but
because it is \emph{definitionally} resistant to the mathematical tools we have.  The
proof, if it exists, will require one of the following:
\begin{enumerate}[label=\upshape(\alph*)]
  \item A direct construction of SURF (W2);
  \item An unconditional proof of GUE gap universality (W3);
  \item A fundamentally new tool that is simultaneously deterministic, upper-bounded,
    and non-positivity-based (bypassing the capstone N1).
\end{enumerate}
This programme has proved that no currently known tool achieves any of these three.

\medskip
\noindent\emph{``Una falsa victoria sería peor que un fracaso.''} --- Programme motto.

%% =====================================================================
\begin{thebibliography}{99}
%% =====================================================================

\bibitem{RP7}
D.~A.\ Trejo~Pizzo,
\emph{A rigorous localized Weil detector},
Preprint (2025).
\bibitem{RP8}
D.~A.\ Trejo~Pizzo,
\emph{Faithful Kreĭn-space realization of the Weil form},
Preprint (2025).
\bibitem{RP9}
D.~A.\ Trejo~Pizzo,
\emph{Stable local arithmetic fingerprints from localized Weil forms},
Preprint (2025).
\bibitem{RP10} D.~A.~Trejo Pizzo, \emph{The localized-Weil anatomy as a complete local invariant of $L$-functions: Euler products, Satake recovery, and functoriality}, preprint, 2026.

\bibitem{RP11}
D.~A.\ Trejo~Pizzo,
\emph{The band-limited Weil--Carleson form},
Preprint (2025).
\bibitem{RP12}
D.~A.\ Trejo~Pizzo,
\emph{A Lyapunov theorem for the de Bruijn--Newman flow},
Preprint (2025).
\bibitem{RP13}
D.~A.\ Trejo~Pizzo,
\emph{Modern routes to RH and the wrong-sign capstone},
Preprint (2025).
\bibitem{RP14}
D.~A.\ Trejo~Pizzo,
\emph{The $\om$-hierarchy, branching random walks, and multiplicative chaos},
Preprint (2025--2026).
\bibitem{RP15}
D.~A.\ Trejo~Pizzo,
\emph{Why the Riemann Hypothesis resists: finitization obstruction},
Preprint (2026).
\bibitem{RP16}
D.~A.\ Trejo~Pizzo,
\emph{Pontryagin index and finite defect},
Preprint (2026).
\bibitem{RP17} D.~A.~Trejo Pizzo, \emph{An unconditional finite bottom for the localized Weil form: a verified second-moment identity, an almost-everywhere clean bound, and the residue as a uniform short-interval prime variance}, preprint, 2026.

\bibitem{RP21}
D.~A.\ Trejo~Pizzo,
\emph{The Lefschetz dichotomy},
Preprint (2026).
\bibitem{RP25} D.~A.~Trejo Pizzo, \emph{The localized Weil form and its Kre\u\i n--Pontryagin index: localization, obstruction, and rigidity}, preprint, 2026.

%% External references

\bibitem{dB88}
L.~de~Branges,
\emph{Hilbert Spaces of Entire Functions},
Prentice--Hall, Englewood Cliffs NJ, 1968.

\bibitem{KS00}
J.~P.\ Keating and N.~C.\ Snaith,
\emph{Random matrix theory and $\zeta(1/2+it)$},
Comm.\ Math.\ Phys.\ \textbf{214} (2000), 57--89.

\bibitem{Ten15}
G.~Tenenbaum,
\emph{Introduction to Analytic and Probabilistic Number Theory},
3rd ed., Grad.\ Stud.\ Math.\ \textbf{163}, AMS, 2015.

\bibitem{Weil52}
A.~Weil,
\emph{Sur les ``formules explicites'' de la th\'eorie des nombres premiers},
Comm.\ S\'em.\ Math.\ Univ.\ Lund (1952), 252--265.

\end{thebibliography}

\end{document}
